% Bản tiếng Việt tương ứng của phần tóm tắt phương pháp.
% Kết quả định lượng được trình bày riêng trong chương thực nghiệm.
Phân đoạn polyp trong video (Video Polyp Segmentation -- VPS) là một
bài toán quan trọng trong lĩnh vực nội soi đại tràng có sự hỗ trợ của
máy tính. Những tiến bộ gần đây của các mô hình nền tảng cho phân đoạn
ảnh, đặc biệt là họ mô hình Segment Anything Model (SAM), đã mở ra
triển vọng trong việc xác định ranh giới tổn thương và lan truyền
kết quả phân đoạn theo thời gian. Tuy nhiên, việc phân đoạn polyp
một cách đáng tin cậy vẫn gặp nhiều thách thức do sự biến đổi về
hình dạng và đặc điểm thị giác của tổn thương, hiện tượng che khuất,
cũng như sự tương đồng giữa polyp và các mô xung quanh. Bên cạnh đó,
hiệu quả của các phương pháp phân đoạn dựa trên prompt phụ thuộc
đáng kể vào độ chính xác của bước khởi tạo đối tượng mục tiêu, trong
khi các khung bao được sinh tự động bởi bộ phát hiện bên ngoài như
YOLO có thể thiếu chính xác hoặc gây nhập nhằng về vị trí.

Để giải quyết những thách thức trên, khóa luận đề xuất một khung
phân đoạn polyp tự động trong video, kết hợp khả năng phát hiện
tổn thương của YOLO, cơ chế phân đoạn theo prompt của SAM~2 và
các cơ chế lan truyền theo thời gian với bộ nhớ hiệu quả của
Lean-SAM2. Khung phương pháp tập trung giải quyết ba vấn đề bổ
trợ trong phân đoạn video: (1) khởi tạo đối tượng mục tiêu đáng
tin cậy, (2) học và tổng hợp các biểu diễn không gian có tính
phân biệt cao, và (3) quản lý bộ nhớ thời gian nhằm duy trì
tính nhất quán của đối tượng mục tiêu. Trên cơ sở đó, ba mô-đun
được đề xuất.

Thứ nhất, mô-đun \textbf{Khởi tạo Prompt có Hướng dẫn bởi Đặc trưng
(Feature-Guided Prompt Initialization -- FGPI)} xây dựng chiến
lược khởi tạo có nhận biết đối tượng mục tiêu bằng cách kết hợp
các khung bao do YOLO tạo ra với những biểu diễn thị giác ở
các tầng đầu của bộ mã hóa ảnh. Trong mỗi vùng được phát hiện,
FGPI xây dựng tập điểm neo giả định của đối tượng có ràng buộc
không gian và đa dạng về đặc trưng, đồng thời khai thác sự tương
phản giữa đặc trưng tiền cảnh và hậu cảnh để lựa chọn các điểm
prompt dương tính và âm tính giàu thông tin. Những điểm prompt
được suy ra từ đặc trưng được kết hợp với khung bao ban đầu và
đưa đồng thời vào bộ mã hóa prompt của SAM~2, hướng tới cải thiện
khả năng xác định đối tượng trong bước phân đoạn khởi tạo.

Thứ hai, mô-đun \textbf{Tổng hợp Đặc trưng Đa giai đoạn có Nhận biết
Vùng (Region-Aware Multistage Feature Aggregation -- RMFA)}
nâng cao chất lượng biểu diễn không gian bằng cách kết hợp các
đặc trưng có độ phân giải cao từ bộ mã hóa ảnh với những biểu
diễn mức sâu đã được điều kiện hóa bởi bộ nhớ thời gian. Thông
qua cơ chế chú ý theo kênh có nhận biết tiền cảnh--hậu cảnh
và khả năng phân biệt đặc trưng theo từng vùng, RMFA khai
thác thông tin không gian và ngữ nghĩa bổ trợ giữa nhiều mức
đặc trưng. Các biểu diễn thu được được sử dụng để dự đoán
thành phần hiệu chỉnh phần dư cho logits phân đoạn, nhằm tinh
chỉnh ranh giới polyp và hạn chế sự can nhiễu từ hậu cảnh,
đồng thời giữ nguyên bộ giải mã phân đoạn đã được huấn luyện
trước.

Thứ ba, mô-đun \textbf{Tinh chỉnh Bộ nhớ Thời gian có Nhận biết Đối
tượng Mục tiêu (Target-Aware Temporal Memory Refinement -- TTMR)}
xây dựng chiến lược lựa chọn bộ nhớ dựa trên đối tượng mục
tiêu phục vụ quá trình lan truyền phân đoạn trong video. Bằng
cách kết hợp độ tương đồng với các điểm neo của đối tượng,
mức độ liên quan của từng token bộ nhớ và tính đa dạng theo
thời gian, TTMR ưu tiên duy trì các biểu diễn lịch sử có khả
năng phân biệt cao, đồng thời giảm thông tin dư thừa trong
giới hạn dung lượng bộ nhớ. Cơ chế này hướng tới bảo toàn
đặc trưng nhận dạng của đối tượng qua các khung hình liên
tiếp và hạn chế sự tích lũy sai số phân đoạn khi đối tượng
thay đổi hình thái, bị che khuất hoặc xuất hiện hiện tượng
trôi dự đoán trong quá trình lan truyền.
